\section{Related Work}
\label{sec:related_work}

Our work connects three research areas: benchmark concentration studies, foundation model impact analysis, and meta-science of machine learning.

\subsection{Benchmark Concentration Studies}

Koch et al.\ (2021) provide the most directly relevant prior work, documenting benchmark concentration patterns from 2015--2020 using Papers With Code and Semantic Scholar data. They found Gini coefficients of 0.6--0.7, demonstrating that research concentrates on fewer datasets over time, with elite institutions dominating influential dataset creation. Their analysis reveals that 15\% of ML datasets account for over 80\% of benchmark usage---a power-law distribution suggesting systemic concentration.

However, Koch et al.'s analysis ends at 2020---precisely when foundation models began transforming the landscape. Their methodology uses static snapshots rather than time-series analysis, precluding detection of structural breaks. Our work extends their temporal coverage through 2024 and introduces change-point detection to identify \emph{when} concentration dynamics shifted.

Raji et al.\ (2021) critiqued benchmark practices from a construct validity perspective, arguing that benchmarks framed as measuring ``general'' AI progress often lack construct validity. They document cases where benchmark saturation (near-ceiling performance) fails to translate to real-world capability, and where leaderboard optimization encourages overfitting to specific test distributions. Their work explains \emph{why} benchmark overuse is problematic but does not quantify concentration dynamics over time.

Schlangen (2021) examined ``benchmark rot''---the phenomenon where benchmarks become less informative as models saturate them. This complements our finding that attention shifted to emergent benchmarks: the transition may partly reflect benchmark lifecycle dynamics rather than pure foundation model effects.

\subsection{Foundation Model Impact Analysis}

Extensive work documents foundation model capabilities. Brown et al.\ (2020) demonstrated GPT-3's few-shot learning, showing that scale enables emergent in-context learning without task-specific fine-tuning; Dosovitskiy et al.\ (2021) showed Vision Transformers matching CNNs on image classification while enabling unified architectures across modalities; Devlin et al.\ (2019) established BERT's transfer learning paradigm that made pre-training + fine-tuning the dominant approach.

Bommasani et al.\ (2021) coined ``foundation models'' and analyzed their systemic implications, arguing that these models represent a paradigm shift affecting the entire ML research ecosystem---not just model architecture but evaluation, deployment, and safety practices. Their analysis is primarily conceptual; our work provides empirical quantification of one such systemic effect.

Wei et al.\ (2022) documented emergent capabilities---abilities that appear only at scale and cannot be predicted from smaller models. This emergence creates demand for new benchmarks: capabilities like chain-of-thought reasoning require evaluation methods that traditional benchmarks (designed for discriminative tasks) cannot provide.

\subsection{Meta-Science of Machine Learning}

Wagstaff (2012) argued that ML research overemphasizes benchmark performance at the expense of scientific understanding. Lipton and Steinhardt (2019) documented troubling trends including failure to distinguish contributions and lack of ablation studies. Our work contributes to this meta-scientific literature by treating the benchmark ecosystem itself as an object of study.

\subsection{Our Position}

We build on Koch et al.'s foundation while addressing three limitations: (1) temporal coverage through 2024, (2) dynamic rather than static analysis via change-point detection, and (3) mechanism verification beyond descriptive statistics. We complement Raji et al.'s validity critique with quantitative dynamics, and ground Bommasani et al.'s conceptual analysis in empirical measurement.
